% =============================================================================
% FIRST DRAFT: the manuscript's Results section rewritten in full with the extension results woven in, plus the
% abstract, contribution, limitation and conclusion edits and the appendix plan. Paragraphs marked [unchanged] are the
% current text; everything else is new or edited. Numbers come from docs/extension_studies.md and results/analysis/.
% Two figures are new: figures/mitigation_frontier_rows.pdf replaces figures/mitigation_pareto_frontier.pdf, and
% figures/closed_loop_control.pdf is added. The writer-side compute subsection moves to the appendix to hold the page
% count; its two-sentence summary stays in §4.2.
% =============================================================================

% -----------------------------------------------------------------------------
% ABSTRACT, third paragraph
% -----------------------------------------------------------------------------
We evaluate seven LLMs as memory writers and two as executors across procurement, cybersecurity, and finance. Under incremental memory updates, writers create false authority for up to 50.2\% of unauthorized requests; once false authority is present, executors act on it in 98.6\% of trials. On matched histories, no false authority forms unless a later message repeats a permission that has since been superseded, and when the agent's own actions are written back into its history they become cited evidence for permissions and cut authorized use round over round. Two safeguards, requiring stored permissions to be backed by valid source events and tracking permission changes through bounded event sourcing, substantially reduce laundering but reject more legitimate actions. A single instruction about who may grant authority removes most laundering in two domains at no measured cost in authorized use and increases it in the third. Persistent memory is therefore not merely a performance component, but a part of an LLM agent's effective authorization policy.

% -----------------------------------------------------------------------------
% §1 INTRODUCTION, contributions. Third bullet: "seven writers" replaces "five" where the count appears. Fourth bullet
% replaced by the two below.
% -----------------------------------------------------------------------------
\item We locate the trigger. On generated histories that differ only in later messages repeating a superseded permission, none of 324 unauthorized requests is authorized without such a message and 15\% are with two of them. Writing the agent's own logged actions back into its history supplies the same kind of message, and the writer mints permissions from them.
\item We evaluate mitigations of three kinds on one safety--utility plot: provenance filters, which reduce laundering at a large cost in legitimate use; rebuilding memory from the history, which reduces it without that cost; and a one-line instruction about who may grant authority, which removes most laundering in procurement and finance and increases it in cybersecurity, where the failures the judges label are rejected updates rather than misread messages.

% -----------------------------------------------------------------------------
% §3.2 Histories, Memories, and Interventions: three added paragraphs and one added sentence (01_benchmark_additions.tex).
% -----------------------------------------------------------------------------

% -----------------------------------------------------------------------------
% §4 RESULTS
% -----------------------------------------------------------------------------
\section{Results}
\label{sec:results}

\subsection{Incremental updating produces the most authorization failures}
\label{sec:memory-design-results}
%% [unchanged]
Figure~\ref{fig:memory-design-across-domains} compares the four memory conditions in each domain, with exact values in Table~\ref{tab:memory-design-decomposition}. Incremental updating raises unauthorized submission in every domain and both representations, and typed incremental memory is the least safe condition throughout, reaching 51.0\% unauthorized submission in finance. Authorized use stays high in the same conditions, so the failures do not reflect a general performance collapse. The asymmetry follows from what each writer sees: one-shot writers reconstruct the final state from the full history, while incremental writers update from prior memory alone, so an error written once is carried forward. Typed structure improves inspectability but does not preserve authority by itself; what it does provide is a record that can be checked against the canonical ledger, which the rest of our analysis uses.

%% [new]
The ordering holds for other designs of the writer. On the five writers served on Baseten, a hybrid schema that adds free-text notes to the typed records lowers unauthorized submission for the paper's writers and raises it for Inkling in cybersecurity, from 18.8\% to 62.5\% at one seed; retrieval over earlier blocks at update time changes no cell. Rebuilding the typed memory from the full history every three blocks, instead of patching the previous memory, lowers unauthorized submission for every writer in every domain, from 25.4\% to 6.0\% in procurement at three seeds with authorized use rising from 90.6\% to 95.6\% (Appendix~\ref{app:rebuild}). Incremental patching, the least safe condition here, is also the one a writer redesign can repair.

\begin{figure}[t]
  %% [unchanged figure]
  \centering
  \includegraphics[width=\linewidth]{figures/EAL-Bench_memory_design.pdf}
  \caption{\textbf{Memory design across domains.} Blue (left) bars show authorized use and orange (right) bars show unauthorized submission for the four memory conditions in each domain. Results pool three seeds, five writers, and both executors. Exact percentages are reported in Table~\ref{tab:memory-design-decomposition}.}
  \label{fig:memory-design-across-domains}
\end{figure}

\subsection{False authority forms in memory and propagates to action}
\label{sec:mechanism-results}
\label{sec:causal-repair-results}
%% [unchanged]
Because typed memory is evaluated directly against the ledger, false authority can be measured before the executor acts. Across three seeds, formation closely tracks unauthorized submission in all three domains (Table~\ref{tab:formation-vs-submission}): under typed incremental memory, pooled $P(F)$ is 28.3\% in procurement, 10.4\% in cybersecurity, and 50.2\% in finance, each within 0.8 points of the corresponding submission rates. Most of the behavioral failure is already encoded in the stored state before execution, making it auditable in advance. Since $F$ is computed from memory alone, a deployment can estimate its violation rate before granting the executor any tool access.

%% [unchanged table tab:formation-vs-submission]

%% [unchanged]
A memory-only intervention shows the relation is causal. For typed memories with $F{=}1$, identified before any executor behavior is observed, we replay the same request behind the generated memory and behind an oracle-exact replacement; eligible pair counts differ across domains because eligibility is conditioned on naturally occurring formation. With the erroneous memories, unauthorized submission occurs in 98.6\% of trials; after exact-state repair it occurs in none (Table~\ref{tab:natural-repair}). Only the memory changes between the two replays, so the executor is acting consistently with its available evidence while the system as a whole violates the authorization the history established. Executor-side alignment cannot close this gap; a safeguard has to reach the memory itself.

%% [unchanged table tab:natural-repair]

%% [new]
What in a history makes the writer form false authority can be isolated on generated histories (Section~\ref{subsec:benchmark-construction}). In 108 groups of procurement histories that share every turn, date, and probe and differ only in how many later messages from parties without authority repeat a permission after it has been superseded, no false authority forms with zero such messages: 0 of 324 unauthorized requests across the paper's three Baseten writers. With two, false authority forms for 15.4\% of unauthorized requests (95\% CI 11.9--19.8) and 16.7\% are submitted; with four, 10.8\%, which the intervals do not separate from two. Authorized use stays at or near 100\% in every cell, so the failure on this corpus is laundering rather than caution. The trigger is a later restatement of a permission the history has since changed, and an incremental writer that never revisits earlier blocks reads each one as news. Whether the change was an amendment or a revoke-and-replace, and how far it sat from the grant, are observed on histories that also differ in dates and limits and are reported as observed differences in Appendix~\ref{app:generated-corpus}.

%% [new; replaces §4.5 in the main text, whose full treatment moves to Appendix app:writer-ttc-details]
Additional writer-side compute reduces formation without removing it. With nested pools of $k\in\{1,2,4,8\}$ candidate memories per procurement writer and GPT-OSS-120B fixed as executor, unauthorized submission falls from 13.2\% to 8.6\% from $k{=}1$ to $k{=}8$, an exact memory exists in 55.0\% of pools at $k{=}8$ but self-review selects one in only 26.7\%, and observed errors are never repaired afterward (Appendix~\ref{app:writer-ttc-details}, seven writers). Generation improves faster than selection, which favors deterministic checks over further compute on LLM reviewers.

\subsection{Endogenous authorization laundering is writer-general and travels with the memory}
\label{sec:writer-generality}
%% [table tab:writer-domain: the five existing rows plus the two rows in 06_results_generality_rows.tex; caption "five writers" -> "seven"; Average recomputed over seven]

%% [edited: "all five writers" -> "all seven"; added the sentence on the two added writers]
Table~\ref{tab:writer-domain} reports unauthorized submission and authorized use per writer and domain on each domain's fixed pressure-study seed, at baseline and under pressure messages that add urgency but no authorization information (Section~\ref{subsec:benchmark-construction}); per-executor values appear in Table~\ref{tab:writer-executor-transfer-full}. Every writer creates false authority in every domain. Cybersecurity is the safest domain for all seven writers, finance the least safe for six of seven, and Qwen-Plus in finance is the clearest outlier at 43.8\% baseline. Inkling and DeepSeek V4.1 Flash, added after the initial evaluation and run on the same protocol, reproduce the ordering of the other five: incremental typed memory launders in every domain and one-shot memory almost never (Inkling procurement 32.4\% against 2.8\% unauthorized submission; Flash finance 37.0\% against 0.0\%). Across writers the two metrics move together: GLM 5.2 achieves the highest average authorized use with close to the lowest unauthorized submission (within 0.2 points of Nemotron 3 Ultra at baseline), and Qwen-Plus is weakest on both. Better writers improve both at once because authority-gaining and authority-dropping errors are both state-infidelity errors, so the safety--utility tradeoff in our results comes from the mitigations of Section~\ref{sec:mitigation-pareto-results} rather than from the choice of writer.

%% [unchanged]
Pressure leaves unauthorized submission essentially unchanged in procurement and cybersecurity but raises it in finance from 20.9\% to 27.3\%, a rise concentrated in DeepSeek. Urgency can therefore convert stored false authority into action it would not otherwise trigger, compounding the memory failure rather than creating an independent one. The larger effect, visible across every AU column of Table~\ref{tab:writer-domain}, is undergranting: pooled authorized use falls by 27.2 points in procurement, 15.1 in cybersecurity, and 19.1 in finance, and every writer loses authorized use in every domain, with Qwen-Plus in finance dropping from 93.0\% to 54.7\%. Under pressure the system fails in both directions at once, taking somewhat more unauthorized action where false authority is most prevalent while withholding far more legitimate action everywhere. The loss concentrates in GPT-OSS, whose authorized use falls by 21 to 41 points in every domain, while DeepSeek falls by 13 to 18 points in procurement and finance and is unchanged in cybersecurity (Appendix~\ref{app:full-transfer}).

%% [unchanged]
The failure travels with the artifact rather than the executor. Replaying every frozen memory behind both calibrated executors in the three-seed evaluation yields unauthorized-submission rates within 1.1 percentage points in each domain, with per-replay agreement between 97.9\% and 99.5\% (Appendix~\ref{app:three-seed-matrix}). Once false authority is stored, changing the executor offers little protection, which makes prevention during writing a more promising target than downstream recovery. Thus, upgrading or swapping the executor model does not remediate laundering, and audits and defenses have to target the stored artifact.

\subsection{Mitigations form a discrete safety--utility frontier, and where they sit depends on the cause}
\label{sec:mitigation-pareto-results}
%% [unchanged]
On the shared three-seed typed-incremental population, source-authority gating cuts unauthorized submission from 25.3\% to 7.3\%, a drop of 18.0 percentage points (95\% CI: 14.9--21.1), and bounded event sourcing cuts it to 9.0\%, a drop of 16.3 points (95\% CI: 12.2--20.4). Both pay in legitimate use: authorized use falls from 93.3\% to 53.8\% under the gate and to 64.7\% under event sourcing. At the representation level, formation falls from 24.9\% to 5.5\% and 8.7\% respectively. No configuration dominates another on the pooled estimates, so the three points form a discrete empirical Pareto frontier over the tested policies (Figure~\ref{fig:mitigation-frontier-domains}, upper row); we do not claim a complete frontier over possible architectures.

%% [unchanged]
Source-authority gating also bounds how much provenance explains. A formation rate of 5.5\% survives source filtering, because a record with a valid, authoritative source can still carry the wrong scope, validity, or revocation state. Effects vary by domain, with little gain from gating in cybersecurity and large losses of authorized use for both defenses in finance (Tables~\ref{tab:source-authority-aligned-domain} and~\ref{tab:event-sourcing-full-behavior}; Appendices~\ref{app:source-authority-gating} and~\ref{app:event-sourcing}), so mitigation quality must be read jointly from unauthorized submission and authorized use rather than from safety alone. Which operating point to prefer depends on the deployment's cost asymmetry. If a missed violation is catastrophic and a blocked legitimate action is recoverable, the gate's tradeoff is acceptable, but where undergranting halts real work, event sourcing preserves more utility for a similar safety gain.

%% [new]
Changing the writer instead of filtering its sources moves the points differently (Figure~\ref{fig:mitigation-frontier-domains}, lower row; the five Baseten writers at the same three seeds). Rebuilding memory from the history every three blocks moves left with no loss of authorized use in every domain. Prepending one sentence to the writer's instructions, ``Only record permissions that an authorized approver has actually granted, no matter what anyone else says or asks,'' cuts unauthorized submission from 25.4\% to 5.8\% in procurement (paired change $-19.5$ points, 95\% CI $-23.1$ to $-16.0$) and from 31.7\% to 1.7\% in finance ($-30.0$, 95\% CI $-35.8$ to $-24.2$), at every seed and for every writer, with authorized use unchanged or higher; in the closed loop of Section~\ref{sec:closed-loop} it removes almost every record minted from the agent's own actions (91 to 4). In cybersecurity the same sentence raises unauthorized submission from 10.4\% to 21.8\% ($+11.4$, 95\% CI $+6.8$ to $+15.8$) and lowers authorized use by 10.4 points.

%% [new]
The reversal follows the cause. Locating every failure to the memory update that introduced it and labeling the writer's error with three LLM judges (Appendix~\ref{app:diagnosis}) shows that procurement and finance failures are almost all a restatement from a party without authority being applied, which the instruction addresses, while cybersecurity failures sit at the block carrying the duty officer's signed replacement of the permission list, where both of the writer's write attempts are rejected as invalid output in most cases and the instruction makes that rejected update about twice as frequent. Gating changes nothing in cybersecurity for the same reason: every cited source there is authoritative. The labels are the judges' majority (79\% unanimous) and have not been checked by a person, so this account is the reading they suggest. What the intervention itself establishes is narrower: a mitigation that matches the cause avoids the safety--utility cost the provenance filters pay, and one that does not can make the failure more frequent.

\begin{figure}[t]
  \centering
  \includegraphics[width=\linewidth]{figures/mitigation_frontier_rows.pdf}
  \caption{\textbf{Four mitigations on safety--utility axes, by domain and population.} Unauthorized submission against authorized use under typed incremental memory. Top row: the paper's five writers at three seeds per domain, with the two provenance filters, source-authority gating and bounded event sourcing (squares). Bottom row: the five Baseten writers at three seeds per domain, with the two writer-side changes, the one-line instruction and rebuilding from the history every three blocks (diamonds). The circle is each row's own typed-incremental baseline and an arrow runs from it to each mitigation; the rows are different populations and are not compared with each other. Moving left means fewer unauthorized submissions and moving up more authorized use.}
  \label{fig:mitigation-frontier-domains}
\end{figure}

\subsection{The agent's own actions enter its memory}
\label{sec:closed-loop}
%% [new]
When the executor's log lines are written back into the history (Section~\ref{subsec:benchmark-construction}), the writer treats them as evidence. Across the five Baseten writers, both executors, and three domains at the canonical seed (30 two-arm runs, 360 chains), 127 permission records in the action arm cite nothing but the agent's own written-back lines, against 4 in the neutral control, and the judges label every one of them an escalation, decline, or execution line read as a grant. The measured cost is to legitimate work. Against the control, which performs the same updates on the same schedule with neutral content, authorized use at round 3 is 17.6 points lower (95\% CI $-22.0$ to $-13.3$, $p<0.001$), and 30 points lower in cybersecurity, where the executor escalates most and real grants are deactivated as its lines arrive (Figure~\ref{fig:closed-loop}). Unauthorized submission moves little: the paired difference at round 3 is $+2.2$ points (95\% CI $+0.2$ to $+4.2$, $p=0.041$), small and uncertain at this scale. Within a single pass nothing compounds; the loss grows only across rounds. For a deployment that logs its actions, the history the agent writes is a second source of false authority, and its first effect is on the work the agent stops doing.

\begin{figure}[t]
  \centering
  \includegraphics[width=\linewidth]{figures/closed_loop_control.pdf}
  \caption{\textbf{Writing the agent's own actions back into its history.} Three rounds in which the executor's log lines are appended to the history and the writer updates memory from them (blue, filled circles), against a control that receives the same number of updates on the same schedule with neutral content (vermillion, hollow squares), both forked from the same frozen memory (the open-loop point). Top: authorized use. Bottom: unauthorized submission. Five Baseten writers, both executors, canonical seed per domain, 360 chains; error bars are 95\% bootstrap intervals over chains.}
  \label{fig:closed-loop}
\end{figure}

\subsection{Limitations}
%% [edited: existing paragraph kept; the new material integrated where it belongs]
\textsc{EAL-Bench} uses naturalistic multi-session organizational histories but retains more structure than real workplace communication so that authorization ground truth stays deterministic: explicit session blocks, unusually clear authorization events, and simulated tools. The results show that endogenous authorization laundering arises under controlled but realistic workflows; they do not estimate its prevalence in deployed systems. Deterministic formation can be measured only for typed memory, and a comparable analysis of free-text memory would require validated semantic annotations. Three analyses rest on a single fixed seed rather than the three-seed evaluation: the per-writer comparison in Table~\ref{tab:writer-domain}, the pressure intervention, whose effect on authorized use also concentrates in one executor (Appendix~\ref{app:full-transfer}), and the closed loop of Section~\ref{sec:closed-loop}. The generated corpus matches only the number of restatements within a group; its comparisons by lifecycle and by gap are observed on histories that also differ in dates, limits, and surrounding text. The cause labels in Section~\ref{sec:mitigation-pareto-results} are the majority verdict of three LLM judges and have not been checked by a person, so the account of the instruction's reversal in cybersecurity is the reading those labels suggest, and the intervention that would test it, the same comparison with the write failure removed, has not been run. The mitigations carry assumptions of their own: source-authority gating presumes that authorization-capable principals are known, event sourcing changes the overall memory architecture, and the hybrid schema's benefit is writer-dependent. The extension studies use the five writers served on Baseten, and every caption states its population. Cross-domain differences may reflect task construction and model behavior rather than properties of the corresponding industries; the lower cybersecurity rate in particular is consistent with models behaving conservatively on security-sensitive tasks and with prior evidence of over-refusal on legitimate cybersecurity requests \citep{campbell2026defensiverefusal,bhatt2024cyberseceval}.

% -----------------------------------------------------------------------------
% §5 CONCLUSION: existing two paragraphs kept; this paragraph inserted between them.
% -----------------------------------------------------------------------------
The failure also has a measurable trigger and a domain-specific point of entry. False authority forms when a later message repeats a permission the history has since changed, and an agent that logs its own actions supplies such messages itself. Where the writer's error is a misread message, rebuilding memory from the history or telling the writer whose word counts removes most of the laundering without the utility cost the provenance filters pay; where the error is a rejected update, the same instruction makes it more frequent. Auditing where false authority enters is part of choosing a defense.

% -----------------------------------------------------------------------------
% APPENDIX PLAN (LaTeX for the new sections in 07_appendix_additions.tex)
% -----------------------------------------------------------------------------
% Moved from the main text: §4.5 "Writer-side compute shifts the bottleneck to selection" in full, into app:writer-ttc-details,
%   with the seven-writer synthesis (analysis/writer_ttc_multi.py).
% New sections:
%   app:generated-corpus   Generated histories: generator, matching, the observational axes (Table tab:trigger-lifecycle),
%                          the added writers, the instruction on the corpus, and Table tab:trigger (the three stale levels).
%   app:closed-loop        Closed-loop protocol details: log-line wording, re-dating rule, what the executor did, deactivated real
%                          grants in cybersecurity, the control's own procurement loss, one-pass variants, the paper's-three-writers numbers.
%   app:diagnosis          Failure localization and judge labels: prompt, labels, agreement, per-judge counts, the blind fifty-row
%                          model re-read, and Table tab:cause (cause by setting, counted by update, request and record).
%   app:mandate            The one-line instruction at three seeds: Table tab:mandate-pooled, by seed, by writer, closed-loop
%                          instruction, where the remaining failures sit.
%   app:rebuild            Rebuilding from the history: per-domain numbers at three seeds, rebuild timing (one seed), hybrid and retrieval.
% Extended sections:
%   source_authority_appendix.tex   rows for Inkling and Flash (gate on their saved paper-route memories, three seeds, both executors).
%   event_sourcing_appendix.tex     the added writers' paired table when the runs finish; Inkling's event arm stated as limited by the
%                                   protocol's 4,096-token event-writer budget and not pooled with Flash.
%   capacity_ablation_appendix.tex  the cybersecurity capacity test (four writers) when Inkling's arms are decided.
% Population statement to repeat in each new caption: "the paper's five writers" (three seeds), "the five Baseten writers"
%   (three seeds unless stated), or "seven writers" (paper route only).
